% !TeX root = ../main.tex
% 总结与展望

\chapter{总结与展望}
\section{总结}
随着嵌入式软件系统复杂度不断增加，传统覆盖率采集与分析方法难以满足多平台、多领域的测试需求。特别是在资源受限的异构嵌入式平台中，缺乏独立文件系统和测试后自动断电等限制，对覆盖率的实时、可靠采集构成了挑战。此外，传统的全量覆盖率检视方法耗时且容易出错。本文在这样的需求下，设计了一个支持异构平台覆盖率采集、能够高效整合和可视化呈现多维度覆盖率数据的系统。

系统的特色主要体现在以下几个方面。通过以下工作，本文为复杂嵌入式平台的覆盖率管理提供了一种新的思路与方法，并为相关领域的研究与实践提供了重要参考。

\paragraph{\textbf{跨平台与异构环境支持}}
系统设计了一套基于共享内存的跨OS通信机制，能够在无文件系统的异构操作系统总，实现实时且可靠的覆盖率数据采集与传输，突破了传统方法在此类环境下的局限。系统兼容LLVM和GCC交叉编译工具链，实现了通用的覆盖率采集、解析、可视化流程，为多平台覆盖率分析提供了可靠的技术支撑。

\paragraph{\textbf{细粒度覆盖率解析与增量计算}}
系统设计并实现了细粒度解析模块，通过模块-文件映射与代码差异分析，精确提取模块级覆盖率数据与增量覆盖率数据，大幅提高了测试的准确性和效率，满足了复杂嵌入式系统的需求。

\paragraph{\textbf{多维度数据融合与可视化}}
系统提供了基于版本、单板、领域、模块、时间等多维度的覆盖率数据融合与可视化方案。通过精细的层次划分和多维度展示，测试人员能够直观地了解覆盖率的各个层面，从而更有效地发现潜在问题，并提升盲点定位效率。

\paragraph{\textbf{自动化任务执行与集成}}
系统将覆盖率采集、解析和可视化任务自动化，并通过 SmartCI 持续集成平台，将任务集成到冒烟测试流水线中，实现了任务的自动触发与分布式执行。该方案显著提高了测试效率，降低了人工干预，同时优化了资源的使用，满足大规模测试的需求。

\paragraph{\textbf{工程可行性与实际应用}}
系统已经在部门内部推广应用，并获得了积极反馈，充分验证了系统在实际工程中的可行性与价值。系统的成功实施为嵌入式软件的多维度覆盖率管理提供了高效、可行的解决方案，具有广泛的应用前景。

\section{展望}

尽管系统在多个方面实现了创新和优化，但仍存在一些限制和可以进一步发展的空间。

首先，系统当前仅支持模块匹配告警和函数未覆盖告警，并且这些告警信息仅在前端应用中进行展示。为了进一步增强系统的监控与提醒功能，后续可以考虑扩展告警类型，例如引入函数覆盖率过低告警、模块覆盖率过低告警等。这些告警能够在实时测试中为开发人员提供更加详细的反馈。此外，系统还可以支持自动将告警信息发送到相关责任人的邮箱，进一步加强问题追踪和督促，确保测试人员能够及时响应并优化测试工作。

其次，当前系统只支持版本级的覆盖率采集与解析，即仅能对主干代码进行覆盖率数据采集与分析。系统可以扩展开发个人级别的覆盖率分析功能，允许对不同开发分支的代码进行独立的覆盖率采集和解析，从而提升团队协作中的代码质量管理。

进一步地，目前用户只能对提交到远程的代码进行覆盖率分析，这是由于SmartCI平台只能拉取远程源代码进行覆盖率采集与解析。为了解决这一问题，后续可以开发拓展脚本，支持用户在本地环境中执行覆盖率采集与解析任务，在本地快速获取反馈，进而避免不必要的提交，减少远程分支中的频繁小规模提交，提高开发效率和代码质量。